\newpage

\ifdefined\useChapters
\chapter{Entre o Sonho e o Relógio}
\else
\chapter{}
\fi

O tempo anda torto desde que voltei. Às vezes acordo sem saber se dormi, às vezes durmo no meio do dia e sonho acordado. Minha irmã diz que é o trauma, que eu devia procurar ajuda. Mas como explicar algo que não parece ter começado nem terminado?

\medskip

Hoje tentei agir normalmente: levantei, escovei os dentes, liguei o chuveiro. Só que a água caía mais devagar do que devia, como se o tempo tivesse preguiça. Quando percebi, fiquei encarando as gotas caindo, uma a uma, até esquecer o que era pra fazer.

— garoto, tá tudo bem aí dentro? — gritou Lúcia, batendo na porta.  
— Tudo. — menti, desligando a água.

\medskip

Há algo que não me larga — o relógio. Desde que acordei, ele não sai do meu bolso. Não lembro de tê-lo pego antes, mas ele estava na mesa do hospital quando despertei, ao lado das flores murchas que Lúcia deixou.  
Um relógio de bolso, antigo, sem marca, o ponteiro dos segundos andando pra trás.

\medskip

Depois do almoço, saí pra caminhar. A rua parecia mais clara do que devia. O ar, leve. O tempo, suspenso.  
Caminhei até a avenida sem destino, só deixando os passos me levarem. Cada sombra, cada vitrine, parecia uma lembrança que o tempo deixou pra trás.  

O mercado do seu Antônio ainda estava lá — o mesmo letreiro gasto, o cheiro de pão velho e café queimado. Entrei mais por impulso do que por vontade. O sino da porta tilintou, e um homem, que eu nunca tinha visto antes, ergueu os olhos do balcão.  

Ele me observava como quem já me conhece. O olhar fixo, quase curioso, e um leve sorriso disfarçado sob a aba de um boné gasto. As mãos escondidas nos bolsos do casaco, como quem evita ser reconhecido.  

Andei pelos corredores estreitos fingindo procurar algo.  
Peguei uma caixa de fósforos, deixei de lado. Toquei num pacote de farinha, sem motivo.  
Sentia o olhar dele me acompanhando, silencioso. Cada passo parecia ecoar no chão de madeira, ritmado, medido — como o tique-taque do relógio no meu bolso.  

O coração batia junto com o som, acelerando.  
Quando cheguei ao fundo da venda, olhei rápido por sobre o ombro e o peguei desviando o rosto, como quem teme ser visto de frente.  

Aproximou-se devagar, voz baixa, tom grave e familiar de um jeito que eu não soube explicar.

— Boa tarde, garoto — disse, sem me tirar os olhos. — Ainda procurando o que perdeu?

— O quê? — perguntei, fingindo leveza, mas o ar pesou.

Ele sorriu, sem humor, como quem testa um segredo antigo.  
— Quando for dormir hoje, leve o seu relógio. Ele vai saber o que fazer.

Virei o rosto por um instante pra pegar o pão e o troco no balcão.  
Quando ergui os olhos de novo, o homem já não estava.  

Seu Antônio surgiu do fundo com o avental amassado, como se nunca tivesse saído de lá.  
— Tá tudo certo aí, garoto?  
— Tava um homem aqui agora há pouco...  
— Que homem? Eu tô sozinho desde cedo.

Olhei em volta. O corredor parecia mais comprido, a luz mais fraca.  
O sino da porta balançava ainda, mas o vento estava parado.  
Apenas o som do relógio no meu bolso preenchia o silêncio.

\medskip

Voltei pra casa tentando fingir que nada disso era importante.  
Lúcia estava deitada no sofá, a TV ligada no mudo.

— Você tá estranho, garoto.  
— Eu sou estranho desde que nasci.  
— Mais do que o normal. Tá com olheiras, fala sozinho.  
— É só insônia.  
— Então dorme.  
— Tô tentando.  
— Então tenta mais.

Ela desligou a TV e subiu pro quarto.

\medskip

Fiquei ali, olhando para ela, tão familiar e ao mesmo tempo distante. Lúcia sempre foi minha confidente, a irmã que conhecia cada pedaço de mim, desde as brincadeiras na infância até os segredos sussurrados no escuro. Mas agora, após o coma, parecia que uma parede invisível se erguera entre nós. O vínculo se tornara tênue, marcado por silêncios e olhares que tentavam alcançar o que o tempo e o silêncio haviam roubado. Ela expressava carinho do jeito dela — com críticas, com reprimendas, com aquele jeito duro que escondia uma preocupação profunda. Eu entendia, sentia isso, mas também sentia falta da intimidade que tínhamos, das risadas fáceis, das conversas que não precisavam de palavras.

— Ei, garoto — ela chamou da porta do quarto, com um sorriso meio torto. — Se continuar assim, vou ter que te internar de novo só pra cuidar de você.

Sorri, mesmo sem vontade, reconhecendo naquela provocação carinhosa um pouco do que ainda nos unia.

\medskip

Fiquei olhando o reflexo do relógio na mesa. O tique-taque parecia mais alto do que o som do mundo.

\medskip

De madrugada, escrevi num papel o que precisava lembrar:  
\emph{dor, medo, foco, retorno.}  
Era o método. A sequência. A tentativa de provocar o mesmo estado do coma.  
Talvez fosse loucura, mas e se não fosse?

Deitei com o relógio na mão. Ele pulsava quente, como se tivesse coração próprio. O som acelerou:  
\emph{tic-tac-tic-tac-tic-tac.}

Fechei os olhos.

\medskip

Um vento subiu pelo chão, atravessou o colchão, o corpo, o peito.  
Por um instante, não senti peso. Nem cama, nem teto.  
A casa sumiu. O tempo também.  
E, no escuro, uma voz que não vinha de fora:

— Você demorou.

O som me atingiu antes do medo.  
Tentei abrir os olhos, mas só via um contorno dourado flutuando diante de mim.  
Era o relógio. Aberto, os ponteiros girando soltos.

— Quem é você? — perguntei, mas a pergunta soou dentro da minha cabeça.

— Você me trouxe. Ou eu te trouxe, ainda não decidi.

O ponteiro parou. Um estalo, depois silêncio.

O chão apareceu. Eu estava no corredor da casa amarela.

As portas, todas fechadas.

Uma delas vermelha.

\medskip

Apertei o relógio contra o peito, tentando ouvir o coração dele.  
O som misturava-se ao meu.  
Tic. Tac.  
Tic. Tac.  
Um som só.  
Um tempo só.

A maçaneta girou sozinha.

\medskip

Talvez eu nunca saiba se aquilo era sonho, lembrança ou castigo.  
Mas de uma coisa eu tenho certeza: foi o relógio que escolheu quando eu voltaria.

E eu voltei.